\chapter{知识的迁移与共享概述}
\label{chap:kts-overview}

上一部分（第8–11章）我们\emph{盯住同一个持续运行的系统}，看它怎样随时间更新、又不丢掉旧能力。本部分\textbf{换一个视角}：不再问“一个系统怎么随时间维护自己”，而是问\textbf{“在一个任务上已经学到的经验，能不能、以及怎样\emph{拿去帮另一个任务}”}，以及\textbf{“多个任务能不能\emph{从头就一起训练}、共享同一套底子}”。

继续用熟悉的花卉系统。设想园区提出了一批\textbf{新要求}：在已经有一个成熟的“物种分类”模型之后，还要\textbf{新增“叶片病害诊断”}（只有很少病害标注）、要\textbf{在照片里“定位花朵位置”}，并且\textbf{物种、病害、定位最好能用同一套图像特征、同一个服务支撑}。这时你面对的就不是“旧系统随时间漂移”，而是一组\textbf{彼此相关、却各不相同的学习问题}。三个自然的问题随之而来：
\begin{itemize}\item \textbf{“认物种”时学到的花瓣纹理、轮廓这些东西，能不能\emph{帮}“看病害”“找位置”？}——这是\textbf{知识的迁移}；\item \textbf{物种、病害、定位能不能\emph{一起}训练、\emph{共享}同一套特征计算，从而省算力、互相补充？}——这是\textbf{知识的共享}（多任务学习）；\item \textbf{能不能让系统经过大量任务的历练，变得“以后遇到\emph{任何}新任务都能\emph{用很少样本快速上手}”？}——这是\textbf{元学习}与少样本适应。\end{itemize}

\textbf{阅读本章你只需要}：任务、模型、损失、训练、泛化、神经网络表示的基本概念。本章是\textbf{整个第三部分的“问题地图”和“概念边界”}，和第8章在第二部分扮演的角色对应：\textbf{先把“跨任务经验利用”这件事的\emph{语言、记号、范式、权限}讲清楚}，具体算法留给后续三章——第~\ref{chap:ket-transfer-adaptation}章讲知识的迁移，第~\ref{chap:ket-sharing}章讲知识的共享，第~\ref{chap:ket-evaluation}章讲怎样严谨评价它们的收益。

\textbf{先把“演进”和“迁移与共享”的分工再钉牢一次}，因为它们常常同时出现：
\begin{itemize}\item \textbf{演进（第二部分）}的关键词是“\textbf{时间}”：同一个系统，新季节、新类别、新事实\emph{陆续}到来，核心矛盾是\textbf{边学新、边保旧、别僵化}。\item \textbf{迁移与共享（本部分）}的关键词是“\textbf{任务之间}”：\emph{来源任务}的经验怎样帮\emph{目标任务}，以及多个\emph{平级}任务怎样共同塑造一个模型。\end{itemize}
两者会交织：第8章“逐月加入新物种”的每一阶段，本质上都在把旧物种经验\emph{迁移}给新物种；但第二部分关心的是“\emph{长期不遗忘}”，本部分关心的是“\emph{这一次跨任务帮助}本身\emph{何时有效、代价多大}”。本章会把迁移侧的机制讲透，至于“迁移之后怎样在长序列里不遗忘”，仍回到第10、11章。

\section{知识迁移与共享的基本概念}
\label{sec:kts-overview-concepts}

\subsection{从“换任务”说起：三类典型场景}

先用同一个花卉系统看清，“跨任务利用经验”到底发生在哪几种情形里，它们\textbf{要利用的东西、要付出的代价都不一样}。

\textbf{场景一：把旧任务的表示，拿去帮一个\emph{新任务}}。认物种时，网络底层早已学会边缘、纹理、花瓣轮廓这些通用特征。新增“病害诊断”时，\textbf{旧的“输出头”肯定不能直接用}：旧的十个物种分数里，根本没有“健康/黑斑病/白粉病”这个答案；但\textbf{中间的特征计算\emph{可能}非常有用}——病斑也是一种纹理和颜色模式。于是最自然的做法是：\textbf{照片先过旧的特征提取器，再\emph{新训练}一个“从特征到病害类别”的预测头}（俗称“冻结骨干、只训头”），或者\textbf{从旧模型参数出发、继续微调}。

\textbf{场景二：让多个任务\emph{共同}训练、共享一套表示}。既然纹理轮廓可能通用，一个更进一步的想法是：\textbf{不让物种、病害各自训一个独立网络，而是让它们\emph{共享同一套特征计算}（共享骨干）、各接各的输出头，损失一起回传}。它有两个潜在好处：\textbf{一是省重复计算}（一套特征喂多个头）；\textbf{二是让不同标签\emph{共同塑造}一个更丰富的表示}。但共享\textbf{不保证对每个任务都合适}——任务之间可能\textbf{\emph{抢}表示}：\textbf{物种分类可能希望表示\emph{忽略}叶面斑点（斑点对认种是噪声），病害识别却希望表示\emph{保留}斑点（斑点正是信号）}，一套共享特征可能两头不讨好。

\textbf{场景三：为“以后的快速适应”提前做准备}。比“帮好某一个新任务”更有野心的目标是：\textbf{让系统经过大量不同任务的训练后，面对一个\emph{从未见过}的新任务，只给\emph{几张}带标签样本就能上手}。这时被“学习”的已经不只是某个预测器，而是\textbf{“快速学会新任务”这种能力本身}，即元学习（12.1.3 展开）。

\textbf{还要识别一种更隐蔽的“换任务”：任务没改名，但\emph{输入数据的来源变了}}。比如同一个“认物种”任务，模型从江南园区搬到塞北园区，\textbf{要预测的类别、评分标准都没变，照片的\emph{分布}却变了}。这通常叫\textbf{领域变化}，它与“任务变化”是\textbf{两个正交的维度}，下一节正式立记号。

\subsection{来源、目标、领域与任务}
\label{sec:kts-overview-source-target}

经验的跨任务利用总是\textbf{\emph{有方向}的}：“从\emph{谁}那里拿经验、拿去帮\emph{谁}”。同一个旧物种模型，拿去帮“换地区认花”可能很有用，拿去帮“叶片病害”就未必。为这种有方向的关系命名：\term{来源}（source，也译作\emph{源}）指\textbf{\emph{提供}已有经验、模型或学习规则的一侧}；\term{目标}（target，也译作\emph{目标域/目标任务}）指\textbf{当前\emph{希望满足}的新要求}。

三个初学者容易忽略的要点：\textbf{其一，来源/目标是“当前请求里的\emph{角色}”，不是某个系统的永久标签}——这次辛苦训出的病害模型，以后完全可以反过来成为“来源”；\textbf{其二，来源\emph{不一定}是一份能下载的数据集}，一个只开放查询的黑盒模型照样是来源（能提供预测输出供蒸馏）；\textbf{其三，方向不能反}：评价迁移好不好，最终看的是\textbf{“\emph{目标}侧”的表现或代价}，而不是来源侧成绩多高（旧地区 99\% 不代表新地区也好）。

描述来源与目标的差异，要拆成\textbf{两个正交维度}：
\begin{itemize}\item \term{领域}（domain）回答“\emph{会观测到什么样的\emph{输入}、它们按什么规律出现}”——江南的花还是塞北的花、光照如何、各类花出现的比例，它管“\emph{输入长什么样、出现得多频繁}”。\item \term{任务}（task）回答“\emph{要预测\emph{什么}、按什么标准\emph{评价}}”——输出物种、病害还是位置，错了怎么扣分，它管“\emph{要做什么}、怎样算做对”。\end{itemize}
两者\textbf{可以独立变化}：\textbf{换地区}主要变领域（任务还是认种）；\textbf{加病害}主要变任务（输入照片可能没怎么变）。一个编码陷阱值得警惕：把类别编码成整数 $1,2,3$ 时，整数本身没有意义，任务必须规定“3 对应哪种物种”，换个数据集同样的“3”可能指完全不同的东西。

\subsubsection{把“在目标上好不好”变成可比较的数：风险与经验风险}

光说“两边不一样”还不够，必须把“迁移到新任务/新领域后还好不好”变成一个\emph{可以比较的数}。先限定在有监督预测问题里，把记号立起来（第一次见这些符号，逐个对齐）：
\begin{itemize}\item 来源、目标的\textbf{预测输入空间}分别是 $X_S,X_T$，\textbf{预测输出空间}是 $Y_S,Y_T$；\item $\mathcal P_S,\mathcal P_T$ 是评价两侧任务时所用“输入与标签”的\textbf{联合分布}；\item $S_S,S_T$ 是实际用来\emph{估计}风险的\textbf{有限评价样本集}。\end{itemize}
对任意一侧 $D\in\{S,T\}$，预测器 $f_D:X_D\to Y_D$ 给每张照片一个答案，损失 $\ell_D(\hat y,y)$ 给预测打分（越小越好，0 通常是全对）。我们真正关心的是“在这一侧\emph{所有}照片上的\emph{平均}表现”，于是按分布取期望，定义\textbf{风险（期望损失）}：
\begin{equation}
  R_D(f_D)=
  \mathbb E_{(\bm{x},y)\sim\mathcal P_D}
  \bigl[\ell_D(f_D(\bm{x}),y)\bigr],
  \qquad D\in\{S,T\}.
  \label{eq:kts-overview-risk}
\end{equation}
\textbf{特别声明}：训练样本如何产生要\emph{另行}指定，\textbf{\emph{不默认}训练样本与评价样本独立同分布}——在迁移场景里这个默认恰恰常不成立（训练在源域、评价在目标域）。

式~\eqref{eq:kts-overview-risk}立刻给出\textbf{三条“能不能直接比较”的纪律}：
\begin{itemize}\item \textbf{只有两侧\emph{接口相同}（输入输出空间对得上）、\emph{标签含义相同}、\emph{损失相同}时，才能用同一个 $f$ 比较 $R_S(f)$ 与 $R_T(f)$}；此时 $R_T(f)>R_S(f)$ 才说明“换到目标域后恶化了”。\item \textbf{\emph{换了任务}，风险之差\emph{不能}当性能变化}：认种与定位的损失计分对象不同（一个罚错类别、一个罚框偏多少像素），两个数单位不同，相减没有意义。\item \textbf{\emph{即使任务没变}，\emph{没有目标标签}也\emph{算不出}目标错误}：只有照片 $\bm{x}$ 没有正确答案 $y$，就无法从照片本身知道预测对不对——这正是无监督领域适应最难的地方。\end{itemize}

现实中你通常只拿到\emph{有限}的几张带标签照片。设目标侧有 $n_T\ge 1$ 个带标签评价样本 $S_T=\{(\bm{x}_i,y_i)\}_{i=1}^{n_T}$，逐张算同一种损失、求和再除以 $n_T$（除以 $n_T$ 是为了不让样本多少改变平均尺度），得到\textbf{经验风险}：
\begin{equation}
  \hat R_T(f_T)=\frac{1}{n_T}
  \sum_{i=1}^{n_T}\ell_T(f_T(\bm{x}_i),y_i).
  \label{eq:kts-overview-empirical-risk}
\end{equation}
帽号“$\hat{\cdot}$”表示这是用\emph{有限样本}算出的\emph{估计}。一条关键纪律（直接关系到为什么要留测试集）：\textbf{若样本独立抽自目标分布、且预测器\emph{固定}（没看过这批样本）}，$\hat R_T$ 可作为 $R_T$ 的估计；\textbf{但若你把这批样本拿去训练或调参、选模型，它们就参与了学习决策}，模型已“见过并迎合”它们，其上的平均损失\emph{偏乐观}，不能再当独立测试结果。最后提醒：这个平均\textbf{没包含}“旧能力保留得好不好、预算超没超”，那些要另列目标或约束。

\subsubsection{一个常用特例：协变量偏移，以及无标签适应的边界}

“同类花卉换地区”时可固定接口不变（$X_S=X_T$、$Y_S=Y_T$，物种含义、损失都一样），只允许联合分布变化。联合分布可拆成两块：\textbf{输入边缘分布 $\mathcal P_D^X$}（照片怎样出现：逆光多不多、某种背景占比多高）与\textbf{条件分布 $\mathcal P_D(y\mid\bm{x})$}（给定这张照片后标签按什么规律决定）。当你\textbf{额外假定}“两域共享给定输入后的标签规律”，即 $\mathcal P_S(y\mid\bm{x})=\mathcal P_T(y\mid\bm{x})$ 不变、\emph{只有}输入边缘分布变了，这种情形称为\term{协变量偏移}（covariate shift，“协变量”即输入 $\bm{x}$）——“同样长相的东西该是什么还是什么，只是各种长相出现的频率变了”，此时可靠\emph{重加权}源样本适配。若连“物种怎么划分/标注标准”都变了，就超出分布变化、需要重新定义任务，重加权也救不了。

最后点一条理论边界，避免盲目乐观：即便只在“固定输入输出空间、仅分布变化”的范围内，\textbf{要给出风险界，仍需要某种“连接源域与目标域”的条件}。直觉是：对源域里\emph{从没见过}的那部分目标区域，就算给\emph{无限多}的\emph{无标签}目标照片，也\emph{可能}无法确定其标签（同样的输入分布可对应两套相反的标注规律，无标签数据分不开它们）。但“不可学习”\emph{不等于}“所有无标签适应都无效”，真正含义是：\textbf{必须再补“两域有联系”（如共享标签规律）或“目标区域被源域覆盖”的额外依据，迁移才有理论保证}。严谨推导见\BookChapterRef{chap:transfer-learning-theory}{卷3《基础、理论与可信性》第9章“迁移学习的可学习性”}。

\subsection{迁移、多任务与元学习}
\label{sec:kts-overview-paradigms}

有了来源/目标和风险语言，就能给三个最常被混用的范式正式下定义。

\textbf{其一是\term{迁移学习}（transfer learning）}：\textbf{目的是“用\emph{来源}的模型或样本，\emph{帮助\emph{目标}}学得更好/更省”}。它\emph{有方向、有终点}（从 A 迁到 B），是否成功\emph{只能}由“目标侧的表现或学习代价”判断，两类收益：\textbf{一是在\emph{相同}目标测试口径下错误\emph{下降}}（比不迁移更准）；\textbf{二是达到\emph{同一}错误容限所需的目标标签\emph{减少}}（更省标注）。\textbf{\emph{仅仅}“用了来源”不是收益证据}（也可能帮倒忙，即负迁移），必须和“不用来源”的对照比。按来源经验与目标的关系，这类“有方向地再用经验”还可细分为两种：\term{模型与样本复用}（model and sample reuse，把已存样本、条目或模型直接用于当前预测或学习、任务不必变，如再读旧照片）与\term{模型迁移}（model transfer，把一个任务上学到的表示/参数拿去帮另一个任务）；而“边用边改、长期保持”的\term{持续模型维护}（continual model maintenance）则是第二部分的主题，这里只把名字分清。

\textbf{其二是\term{多任务学习}（multi-task learning, MTL）}：\textbf{考察“\emph{多个}目标怎样通过\emph{共同学习、共享}表示或参数}”。它没有单一的“来源到目标”方向，而是多个任务\textbf{平起平坐、互相塑造共享部分}——物种损失和病害损失一起回传、共同改变同一个表示。关键判据是“它们的训练要求是否\emph{作用于同一个共同学习对象}（共享骨干/参数）”。一句话辨析：\textbf{迁移是“从旧的帮新的”（有先后、有方向）；多任务是“几个一起学、共用一套底子”（同时、对称）}。二者可叠加，但概念要先分清。

\textbf{其三是\term{元学习}（meta-learning，字面“关于学习的学习”）}：如果准备阶段的目的\emph{不只是}训好当前一个预测器，而是让\emph{以后遇到的任何新任务都更容易学会}，就进入元学习。普通学习是用照片学会认花；元学习则是用\textbf{一大堆不同任务的学习经验}，学会“怎样快速学会一个新任务”。历史任务经验可以在三个不同层面做这种“准备”：准备\term{模型参数先验}（parameter prior，即在看到新数据前对参数取值的一个概率分布，表达历史更偏好哪类解）、准备一条好的\textbf{学习规则/更新规则}（如 12.3 历史讲的学习优化器、MAML 的初始点），或准备一套可直接检索的\textbf{参照案例}（少样本按相似度对照）。

这带来一个\textbf{抽样单位的根本变化}（理解元学习的钥匙）：\textbf{普通学习抽样的是“照片”；元学习抽样的是“学习问题（任务）}”。记号上：$\mathcal T_k$ 表示第 $k$ 个任务（每个任务自带输入/输出含义、损失、任务内分布，是一个完整小题）；$\Pi$ 是任务的抽样分布（这些小题从哪个“题库”出），$\mathcal T_k\sim\Pi$ 表示先抽一个任务、再从该任务里取照片。于是“训练”变成两级：先抽很多任务，每个任务里再给少量样本。

用花卉把元学习的\textbf{“一集（episode）”}走一遍：\textbf{先抽任务}——这次抽到“区分月季与百合”，下次抽到“区分菊花与兰花”；\textbf{再从抽到的任务里取两类照片}：
\begin{itemize}\item 少量\textbf{带标签的“\emph{支持}照片”（support set）}：用来\emph{当场适应}（看着“这是月季、那是百合”学会区分）；\item \textbf{另取的“\emph{查询}照片”（query set）}：用来\emph{考适应之后还认不认得}（标签不给模型看，只用来打分）。\end{itemize}
评价顺序不能乱：\textbf{先在每个任务内部用查询照片汇总“适应后”的损失，再\emph{跨}很多任务\emph{平均}}——这个跨任务平均衡量的才是“你这套‘学习准备’好不好”。两条纪律随之而来：\textbf{在已见任务上拟合好，不等于会用少量经验适应新任务}（可能只是背熟旧题），所以评价必须在\emph{没参与训练的新任务}上、只给少量样本；\textbf{对不属于 $\Pi$ 覆盖范围的目标不能自动继承结论}（用“区分各种花卉”的题库练出的快速适应，不保证对“识别汽车型号”同样有效）。一个常见“伪元学习”对照：把所有任务照片混成一个大样本集训一个固定分类器，只检验了混合分类准确率，\emph{完全没检验}“全新任务只给几张支持照片时怎么当场学会”。

\subsection{信息条件：你到底能读到什么}
\label{sec:kts-overview-information}

\textbf{两个都叫“迁移学习”的方案，实际能拿到的样本、模型、反馈可能天差地别}，所以比较方法前必须先盘点访问权限，而不能只看方法名。

\textbf{来源侧}从信息全到不全大致有几档：\textbf{源照片（原始样本）可读}（能检查覆盖、回放、重加权，最常因隐私拿不到）；\textbf{源模型参数可读可写}（能复制、冻结、局部修改）；\textbf{只能查询黑盒预测接口}（只拿到\term{预测行为}（predictive behavior），即旧模型返回的类别/概率，没有中间特征、也不能对源参数反传，想蒸馏只能靠“问—答”）。还要把“取得\emph{可执行模块}（能跑起来）”和“读到它的\emph{参数数值}”分开——前者可能只是个加密可调用件。一个重要澄清：\textbf{源样本不开放，排除的只是“直接使用原始样本”，并没有排除“模型已保存在参数/输出里的经验”}——即使看不到原始照片，训练好的源模型或其概率输出仍携带可迁移信息（这正是黑盒蒸馏可行的原因）。

\textbf{目标侧}也要分四类信息：\textbf{目标标签}（最宝贵，通常很少，能直接算监督损失）；\textbf{无标签目标输入}（能估计目标输入长什么样、帮你对齐输入分布，但\emph{不直接给}真实标签）；\textbf{任务描述}（“要识别叶片病害”这几句话，说明判断什么，却不等于逐图标注）；\textbf{反馈时机}（鉴定结果若在预测之后才返回，只能支持以后的更新、不能用于此前那次预测）。由此有一条实验公平性纪律：比较方案时，既要对齐“可用样本数量、模型访问范围”，也要对齐“这些信息\emph{何时}可用”。

\textbf{有了权限清单，就能给几个名字相近、权限相反的常用设定划清边界}：
\begin{itemize}\item \term{领域自适应}（domain adaptation，DA）：\textbf{\emph{允许}在学习阶段\emph{使用目标域样本}}（通常无标签、也可能少量有标签），研究怎样借助来源改善目标域预测；其中无标签 DA、半监督 DA 是不同条件，不能用一个名字混用。\item \term{领域泛化}（domain generalization，DG）：\textbf{权限恰好相反——训练和模型选择阶段\emph{完全不用}目标域样本}，把在多个源域学好的模型直接拿去用在一个\emph{未见}领域，赌的是学到的规律足够通用。\end{itemize}
要防一个误判：\textbf{预测时模型当然要读当前这张照片}（否则没法预测），这是执行预测所必需，\textbf{不等于“被允许用目标样本重新调参}”。必须把“预测执行（读当前输入出答案）”和“按目标样本调整模型（更新参数或状态）”分开记录——一个号称 DG 的方法若上线后偷用一批目标数据重新调参，就跨入了下面的测试时适应。

\term{测试时适应}（test-time adaptation, TTA）允许系统在已上线使用阶段，根据当前到来的目标输入改变自己的参数、统计量或其他保留状态，通常只用无标签目标、靠自监督损失更新，来一张适应一张。TTA 协议必须写死：能不能看到完整批次、能不能跨样本保存状态、状态何时重置、有无反馈权限；一旦按目标样本做了调整，就不能再当标准领域泛化报告。

最后，访问权限之外还有授权、隐私、恶意来源三道关，分别按\BookChapterRef{chap:trust-privacy}{卷3《基础、理论与可信性》第16章“隐私保护”}和\BookChapterRef{chap:trust-security}{卷3《基础、理论与可信性》第17章“对抗安全”}检查。一个反直觉事实：就算“训练样本不离开采集地”或“只开放预测接口（参数不外发）”，模型/接口仍可能泄露训练样本的\emph{成员身份}（某人是否在训练集中）或敏感属性（成员推断、属性推断攻击），所以“数据不出域/只给黑盒”是有用的工程措施，却不能单独作为隐私已保证的证明。

\section{知识迁移与共享的研究内容}
\label{sec:kts-overview-research}

把“目标、来源、权限”描述清楚后，本部分剩下三章各司其职，研究内容可以概括成“\textbf{先准备、再适应、共训、最后评价}”。

\textbf{第一块：来源经验怎样\emph{支持}一个具体新目标、需要\emph{多少代价}}，由第~\ref{chap:ket-transfer-adaptation}章（知识的迁移）负责。它关注两件事：\textbf{其一，来源经验\emph{怎样支持}目标语义}（旧表示里哪些东西帮得上新任务、要冻结哪些层、微调哪些层、怎么换或新增输出头）；\textbf{其二，\emph{形成}这个能力需要多少学习代价}（多少标签、多少步、多少算力）——迁移的价值正体现在“用更小代价达到同等能力”。少样本、零样本、参数高效适应（适配器、LoRA）、面向新输出语义（检测、分割）的改造，都在这一章。

\textbf{第二块：多个任务怎样\emph{共同}训练、协调彼此要求}，由第~\ref{chap:ket-sharing}章（知识的共享）负责。它研究三个共同学习特有的设计问题：\textbf{共享哪些计算}（哪几层骨干共用、哪部分各任务私有）；\textbf{如何协调多个任务的损失}（量纲不同会互相压制，需要加权/平衡，如物种损失与病害损失怎么配比）；\textbf{如何分配更新机会}（这批数据更新哪个任务、采样比例如何）。

初学者常把“多任务共享”和第6章“知识融合”混为一谈，这里给一个清晰分界：\textbf{如果两个模型\emph{已经分别学好了}（两个现成模型在手），那是第~\ref{chap:knowledge-fusion}章的地盘}（集成、参数融合、蒸馏）。分界\emph{不在于“有没有训练”}（融合也可以训练一个组合器或学生模型），\textbf{真正分界是出发点}：\textbf{融合}从“已经存在的、现成的多个来源能力”出发（先有成品再合并）；\textbf{共享/多任务}从“多个目标共同训练一个模型的过程”出发（一开始就一起学，没有两个先学好的成品）。

\textbf{第三块：怎样\emph{严谨评价}迁移、共享、融合到底有没有收益}，由第~\ref{chap:ket-evaluation}章负责。迁移是否有益，要在相同目标证据和预算下比较目标成绩、或比较达到同一要求的全部代价；共享是否有益，要\emph{逐个任务}保留收益、受损与不确定性（不能用漂亮的总平均掩盖某个任务的严重退化）；融合还要检查来源能力是否保留、没见过的属性组合能否正确拼到同一对象上。对照设计、基准、重复实验与统计区间，都在这一章。

为了快速定位，表~\ref{tab:kts-overview-chapters}把本部分章节组织成一张地图。用表时注意：\textbf{这张地图不规定处理顺序}（值不值得迁移往往要先做一点有限试训才知道，真实流程在“判断—试训—再判断”间来回）；但\textbf{“评价标准”必须最先明确}，并在全过程保留样本、权限、版本、成本四份清单，否则事后无法复现、无法证明收益。

\begin{table*}[htbp]
  \centering
  \small
  \begin{tabularx}{\linewidth}{@{}p{50mm}XX@{}}
    \toprule
    章节 & 核心研究问题 & 主体节的组织依据 \\
    \midrule
    第~\ref{chap:kts-overview}章：知识的迁移与共享概述 & 跨任务经验利用提出什么问题 & 概念、范式、信息条件与研究地图 \\
    第~\ref{chap:ket-transfer-adaptation}章：知识的迁移 & 来源经验如何支持新目标及减少学习代价 & 目标确定、来源准备与目标适应 \\
    第~\ref{chap:ket-sharing}章：知识的共享 & 多个任务如何在共同学习中协调要求 & 任务组织、共享结构、联合目标与资源分配 \\
    第~\ref{chap:ket-evaluation}章：知识迁移与共享的评价 & 哪些证据说明预期收益已满足 & 指标、样本与基准、协议和证据分析 \\
    \bottomrule
  \end{tabularx}
  \caption{第三部分章节地图。中列是要回答的问题，右列是各章的拆解方式；同一请求可同时涉及多章，基础机制在主要归属处完整解释，其他章节调用并补充当前任务的条件。}
  \label{tab:kts-overview-chapters}
\end{table*}

还要强调一个\textbf{贯穿全部分的反馈回路}：共享表示是把双刃剑——它可能恰好保留了某个新任务需要的线索，也可能为了分类而舍弃了这些线索；而一旦你为某个任务改动了共享部分，就必须\emph{重新测试}所有共用它的任务。所以评价总原则是：\textbf{不能只留下“某一个环节的最佳数值”，必须同时检查三样——新目标达成没有、相关旧/兄弟任务保住没有、代价花了多少}。任何只报单点最优、不报其余任务和成本的结论，都是不完整的。

\section{知识迁移与共享的研究历史}
\label{sec:kts-overview-history}

跨任务利用经验的研究，历史上主要沿三条\emph{并行}线索展开：\textbf{其一是“任务之间\emph{有没有可共享的规律}、网络的哪一层可迁}”；\textbf{其二是“怎样为\emph{以后的学习}提前做准备，让少量样本就能快速适应}”；\textbf{其三是“多个\emph{各自学好}的能力怎样融合}”。读图/读史前立三条规矩：线索长期并行、互相交织，不是先后接力；选某篇论文定位某个思想，不等于说它是唯一来源或“首次”；后来的方法也不意味着替代了前面全部方法。

\subsection{共享规律与可迁移表示的研究}
\label{sec:kts-overview-history-representation}

\textbf{第一条线索：寻找“任务之间可共享的规律”}。最早被系统研究的入口，正是多任务学习。当病害标注不足时，“物种标签”能不能顺带帮同一套表示学得更好？Caruana 在 1997 年的多任务学习研究中系统讲清：\textbf{相关任务的训练信号作用于同一段共享计算，使一个任务能影响另一个任务}——一个任务的样本相当于在帮共享部分学到对另一任务也有用的规律\scite{caruana1997}。但要记住收益来自“具体的任务联系”而非“输出头变多”这个形式：花卉里两类标签可能共同要求保留叶片细节（互相帮忙），也可能需要不同细节（互相干扰）；而且多任务学习不同于把几个已独立训练好的模型简单并置（那是融合/集成）。

顺着这条线自然会问：共同训练除了做好当前几个任务，能不能顺带为“还没见过的任务”选好一个学习空间（假设空间/表示）？Baxter 在 2000 年把这种“准备”形式化为“\emph{归纳偏置}学习”。这里先解释\term{归纳偏置}（inductive bias）：它指学习者在“众多都能解释数据的可能假设”之间作选择的倾向——即“天然偏好哪一类特征或函数”（偏好平滑的、稀疏的、或分层卷积结构的解）；没有任何归纳偏置，有限数据无法在无穷多个解释间选择。多个任务的共同经验恰好能提供依据，帮你选出“对这一族任务更合适”的偏置。但这种准备有效有条件：任务确实来自某种共同环境、候选假设空间族不能无限复杂；还要分清“在当前这组任务上拟合好”与“为同环境中新任务准备好空间”是两种不同的“好”，前者不能替代后者\scite{ket-baxter2000-bias}。

一旦换到新地区（分布变了），“能共享”的依据就必须接受分布条件的严格检验。2010 年关于领域自适应“\emph{不可学习性}”的理论研究逐条分析：光凭“输入分布之间的关系”“两域都存在一个低误差分类器”“标签机制的某些假定”，够不够支撑学习\scite{bendavid2010impossibility}？核心反例与 12.1.2 呼应：即使目标域无标签输入分布与源域一模一样，也仍可能对应两套相反的标签规律，算法没有任何依据区分；但理论也给出对称的另一面——输入分布变了也不直接证明来源无用（分布不同但条件标签规律可能稳定共享）。两头边界合起来要求你\emph{明确说出到底保留了哪一条预测联系}，既不能只凭“输入看起来很像”当迁移证据，也不能只因“分布不同”就放弃迁移。

到了深度学习时代，问题被细化到“网络的哪一层、用什么操作迁移”。Yosinski 等人 2014 年的经典实验专门把“层”和“迁移操作”当作检查对象\scite{yosinski2014}，发现两件影响深远的事：
\begin{itemize}\item \textbf{不同层有不同的“任务特异性”}：越底层越偏向通用的边缘/纹理特征（容易迁移），越顶层越偏向具体任务的判别（迁移收益下降、甚至有害）；\item \textbf{\emph{切开}原网络会打断层间“共同训练形成的配合}”：即便保留同样的层，只要它们不再和原来的伙伴协同训练，性能也会掉——可迁移性不只是某一层的属性，还和层间耦合有关。\end{itemize}
方法论结论：\textbf{同一层，“冻结后接新头”和“允许继续微调”检验的不是同一种可用性}（前者问特征够不够通用，后者问能不能当好起点）；而且这是特定网络/数据上的结论，不能推出“所有架构都有固定数量的通用层”，必须按层、按目标、按使用方式具体检查。把第一条线索串起来，判断对象沿“任务间有没有联系（Caruana）$\to$ 能选择的表示/假设空间（Baxter）$\to$ 分布变化下可否学习（Ben-David）$\to$ 网络里实际取用哪一层（Yosinski）”一步步变具体。

\subsection{快速适应与学习过程复用的研究}
\label{sec:kts-overview-history-adaptation}

\textbf{第二条线索：让“少量新经验”就能快速适应}。旧表示虽能拿去训病害，却不说明只给几张新标注就够学好。当目标经验非常少时，要求“准备阶段尽量减轻后来学习的负担”，而“准备”至少作用在三个不同对象上，回答三个不同问题：\textbf{准备“怎么\emph{比较}”}（一种样本/类别间的度量，少量样本即可判别）、\textbf{准备“从\emph{哪里}开始}”（一个适合被少量样本微调的参数起点）、\textbf{准备“怎么\emph{修改}}”（一条学出来的更新规则，使几步更新就有效）。下面三个代表工作恰好各占一个对象。

对应“怎么修改”：Andrychowicz 等人 2016 年的“学习优化器（learning to optimize）”把\textbf{更新规则本身当成要学的对象}\scite{ket-andrychowicz2016-learning}。普通训练用人手设计的优化器（如 Adam）更新模型；他们用一个由参数刻画的“更新器网络”，从大量优化问题中学会“看到梯度等训练状态量后应让参数变化多少”，再把这套规则迁到新问题。留下来的是一套更新机制而非某目标模型的最终参数，其有效范围必须按新问题的网络规模、损失形态、优化时长重新检查，不能打包票对任意优化任务有效。

对应“从哪里开始”：2017 年的\term{模型无关元学习}（model-agnostic meta-learning, MAML）选择\emph{参数起点}作为准备对象，也把 12.1.3 那套“一集”机制真正做成算法\scite{finn2017}。准备阶段对每个抽到的任务：先用少量支持样本做几步梯度更新（模拟快速适应），再用同一任务另取的查询样本计算“适应之后”的损失，并用这个损失反过来更新初始参数。角色分工要记牢：\textbf{支持样本负责任务内适应，查询样本负责衡量这种适应对没见过的新样本有没有用}。MAML 最本质的一点是优化“更新之后在新样本上的损失”，而不只是“初始点在混合数据上的损失”——它要找一个“对几步微调特别敏感、一调就好”的起点，而非一开始就最准的点。

对应“怎么比较”：同年的\term{原型网络}（prototypical networks）选择“比较方式”，几乎不学每类的分类器\scite{ket-snell2017-prototypical}。它学一个特征映射，用目标任务里已标注的支持样本算出\emph{每个类别的特征均值（原型）}，新样本按“离哪个原型近”分类，新任务只要每类几个支持样本即可，非常契合少样本。MAML 把目标端参数更新放进准备目标（学好初值）；原型网络则让准备好的表示支持用少量样本直接估计类别位置，适应时基本不更新网络——两者准备目标不同，各有适用范围。

当预训练模型变得非常大，一个现实问题凸显：每个任务都全量微调、再各存一整份模型，参数与存储无法承受。Houlsby 等人 2019 年在自然语言处理中使用\term{适配器}（adapter）：把庞大原网络冻结，只在其中插入一些小型可训练模块，让目标任务只改变这些小模块，每适配一个任务只需训练和保存很小一块增量\scite{ket-houlsby2019-adapters}。这就是\term{模型参数高效适应}（parameter-efficient adaptation，PEFT 是其大家族名）。但要破除一个想当然：\textbf{“可训练参数变少”不能自动推出“需要的标签更少、训练更快、或推理更快}”——适配器仍挂在网络里、推理时照样参与前向（可能还略增开销），训练省的是更新/存储量，各项成本必须分别计算。

2020 年 GPT-3 的少样本报告展示了一种\emph{干脆不改参数}的使用方式\scite{brown2020}：任务说明和几个示例直接以文本进入上下文（提示词），模型只做一次前向就给答案，全程不做梯度更新。它把“提供任务证据”和“训练模型参数”彻底分开：示例确实能改变当次预测，但不构成永久写入——换个会话、撤掉示例，这种行为不保证延续（正是第8章讲的“当次条件化”）。当一个来源在极广泛样本上预训练、再被许许多多下游任务适应时，本书称之为\term{基础模型}（foundation model）：它是“多种目标共享的预训练起点”，不是“已保证解决全部目标的成品”——预训练给的是可适应的底子，每个下游仍要付适应代价、仍可能失败。

\subsection{多来源能力融合的研究}
\label{sec:kts-overview-history-integration}

\textbf{第三条线索：多个各自学好的能力怎样融合}。两个地区模型都已训好，最直接的融合是在预测之后组合它们的判断，即\term{预测集成}（prediction ensemble）：保留每个来源的完整计算，运行时各自预测再按规则组合（类别含义一致时常平均概率或投票）。集成何时真有用？靠的是\textbf{错误互补}：一个犯错时另一个往往对。若两个模型总对同一种背景犯同一个错（都被同样的伪相关骗过），增加来源数并没有补充独立证据，收益就有限。还有个常被忽略的代价：集成要在部署时同时运行并存储所有成员（$N$ 个模型 $N$ 份算力/内存），所以不能只比组合后准不准，还必须算运行与存储成本——这直接引出“想只部署一个模型”的蒸馏路线。

如果部署预算只允许跑一个模型，就要让一个“学生”模型去学习“教师（们）”的行为。Hinton、Vinyals 与 Dean 2015 年的蒸馏工作是代表\scite{hinton2015distilling}：除了硬标签（真实类别），还用\emph{软化}的类别概率（调高温度得到的更平滑分布）作指导，让“非最大概率类别之间的相对倾向”也进入监督——教师不仅说“这是马”，还说“它比像汽车更像驴、比像苹果更像猫”，这些类间相似性是硬标签没有的暗知识（dark knowledge）。能力边界要清楚：学生不需与教师共享参数坐标（结构可不同，学的是输入输出关系），但最终得到多少能力取决于查询输入的覆盖、学习算法和学生容量；它确定的收益是减少部署载体（多个/大教师压成一个小模型），却不保证把教师每一项能力都完整转入。这里按所引 arXiv 版本公开年份 2015 定位，不宣称模型压缩始于该文。

集成融的是“输出”、蒸馏融的是“行为”，Alet 等人 2018 年的模块化元学习（modular meta-learning）则把融合对象推进到“中间计算（模块）”：准备阶段在一系列相关任务上学出一组可复用模块，面对新任务时在允许的结构里选择并组合它们，通过重用模块快速适应\scite{ket-alet2018-modular}。两个条件不能省：接口与“组合语法”必须先规定好；选哪种组合要靠目标证据决定。还要注意最深刻的一个坑——就算每个模块单独都好用，也仍缺“组合起来正确”的证据：假设系统能分别识别“白色”“锯齿叶缘”“位置”，新请求是“找出那朵具有锯齿叶缘的白色花朵”，若图里白花和锯齿叶缘长在两株不同植物上，系统必须把两个属性绑定到同一对象，不能因为“分别检测到白色、也检测到锯齿”就宣布找到目标（这叫绑定/属性归属问题）。这类“没见过的属性/指令组合”是单模块测试根本没检查的关系性能力，必须针对组合结果另设评价，绝不能把单项准确率相加当作组合成绩。

Model Soups（2022）给出另一条“形成单模型”的路径，融合发生在参数层\scite{ket-wortsman2022-soups}：从同一个预训练起点出发，用不同超参/数据顺序得到多个微调结果，直接平均它们的参数得到一碗“模型汤”，常能在不增加推理成本下提升鲁棒性。使用条件要分清：之所以能平均，是因为这些解来自共同起点、处在相容的参数盆地；论文不支持“任意独立训练的模型都能拿来平均”（参数坐标不同义时平均无意义，呼应 12.1 与传感器“维度对齐不等于语义对齐”）。它与预测集成本质不同：权重平均发生在前向计算之前（先合成一套参数再跑一次），预测输出平均发生在各模型各自计算之后（跑 $N$ 次再合答案），两者一般不等价、成本也差很多。

当适配普遍采用适配器后，又出现更省的保存形式：来源不必存整个模型，只存“相对于共同底座的那一点参数增量”。低秩适配器（LoRA 一类）用两个较小矩阵的乘积表示这种增量，每个任务只存一对小矩阵。但融合难题随之而来：不同来源的适配器是各自在不同方向上更新的，它们的“方向坐标”怎样对应？直接把同形状小矩阵相加，方向未必一致。2025 年的 KnOTS 研究先用奇异值分解（SVD）把不同低秩更新变换到一个共同对齐的坐标空间，让方向可比较，然后再应用已有融合规则合并\scite{ket-stoica2025knots}。要准确理解：这是“为融合建立共同坐标”的工程手段，不是声称“任意同形状矩阵相加就能保留来源能力”，收益仍依赖具体模型和基准。

把第三条线索串起来，融合对象被一步步分层、越分越细：预测输出（集成）$\to$ 中间模块（模块化组合）$\to$ 完整模型参数（蒸馏承接、Soups 平均）$\to$ 相对底座的参数增量（KnOTS 融合适配器）。对象不同，需要的访问权限和成本也不同：走“共同运行（集成）”必须在预测时持续访问每一个来源；走“承接模型（蒸馏/单模型）”必须在准备时拿到足够的行为、参数或表示，并测量“转入之后损失了多少能力”。清醒的总结论是：能拿到的模型越来越多，只是扩大了候选集合，并没有免除你做三类评价——冲突（来源互相矛盾怎么办）、来源保留（融合后各来源能力还在不在）、未见组合（拼出新组合对不对）。

图~\ref{fig:kts-overview-history}把三条线索画在同一根等距年份轴上。推荐两个方向交替看：沿横带看每条线索内部“被研究的对象”怎样随时间变具体（如融合线：输出、模块、参数、参数增量）；跨横带看虚线标示的共同问题——“共享表示”和“快速适应”都依赖“任务之间存在可利用规律”，“参数维护”和“多来源融合”都要检查“改动之后还保留了哪些能力”（前者在第二部分、后者在本章）。跨带虚线只标示共同问题，不提供论文间直接师承的证据；横带也不表示某条路线已经终结。

\begin{figure*}[htbp]
  \centering
  % Equal year scale: x = 36 + 4 * (year - 1997), in mm.
% Event years are publication years except the explicitly dated arXiv distillation.
% Bands show concurrent research questions, not their starting/ending dates.
% Dashed arrows link questions, not claims of historical influence.
\begin{tikzpicture}[
  x=1mm,y=1mm,
  font=\figurefont\footnotesize,
  every node/.style={text=BookInk},
  band title/.style={
    align=center,text width=27mm,minimum height=10mm,
    inner sep=0.5mm,font=\figurefont\small
  },
  event/.style={align=center,inner sep=0.7mm},
  leader/.style={draw=BookMuted,line width=0.6pt},
  relation/.style={
    draw=BookTeal,line width=0.8pt,dashed,
    -{Latex[length=1.8mm,width=1.3mm]}
  }
]
  \foreach \level in {0,-32,-64} {
    \fill[BookPaper] (33,\level-12) rectangle (162,\level+12);
    \draw[BookRule,line width=0.6pt] (35,\level) -- (158,\level);
  }
  \draw[BookMuted,line width=0.7pt] (36,19) -- (150,19);
  \foreach \year in {1997,2000,2005,2010,2015,2020,2025} {
    \draw[BookMuted,line width=0.6pt]
      ({36+4*(\year-1997)},18) -- ({36+4*(\year-1997)},20);
    \node[anchor=south] at ({36+4*(\year-1997)},20) {\year};
  }
  \node[anchor=east,text=BookMuted] at (29,21) {年份};

  \node[band title] (representation) at (14,0) {共享规律\\与可迁移表示};
  \node[band title] (adaptation) at (14,-32) {快速适应\\与学习过程复用};
  \node[band title] (integration) at (14,-64) {多来源\\能力融合};

  \draw[relation] (representation.south) -- (adaptation.north);
  \node[anchor=west,text=BookTeal] at (15,-16) {任务规律};

  \foreach \name/\year/\level/\ink in {
    caruana/1997/0/BookBlue,
    baxter/2000/0/BookBlue,
    bendavid/2010/0/BookBlue,
    yosinski/2014/0/BookBlue,
    maml/2017/-32/BookTeal,
    adapters/2019/-32/BookTeal,
    gpt/2020/-32/BookTeal,
    distillation/2015/-64/BookBlue,
    modular/2018/-64/BookBlue,
    soups/2022/-64/BookBlue,
    knots/2025/-64/BookBlue%
  } {
    \coordinate (\name) at ({36+4*(\year-1997)},\level);
    \fill[\ink] (\name) circle[radius=1mm];
  }

  \node[event,anchor=south] (caruana label) at (47,4) {多任务学习\\Caruana（1997）};
  \draw[leader] (caruana) -- (36,2) -| (caruana label.south);
  \node[event,anchor=north] (baxter label) at (67,-4) {偏置学习\\Baxter（2000）};
  \draw[leader] (baxter) -- (49.2,-2) -| (baxter label.north);
  \node[event,anchor=south] (bendavid label) at (96,4) {不可学习性\\2010};
  \draw[leader] (bendavid) -- (88,2) -| (bendavid label.south);
  \node[event,anchor=north] (yosinski label) at (124,-4) {特征可迁移性\\Yosinski等（2014）};
  \draw[leader] (yosinski) -- (104,-2) -| (yosinski label.north);

  \node[event,anchor=south] (maml label) at (108,-28) {MAML与原型网络\\2017};
  \draw[leader] (maml) -- (124,-30) -| (maml label.south);
  \node[event,anchor=north] (adapters label) at (114,-36) {适配器\\2019};
  \draw[leader] (adapters) -- (124,-34) -| (adapters label.north);
  \node[event,anchor=south] (gpt label) at (149,-28) {GPT-3上下文\\2020};
  \draw[leader] (gpt) -- (137.2,-30) -| (gpt label.south);

  \node[event,anchor=south] (distillation label) at (96,-60) {知识蒸馏\\2015};
  \draw[leader] (distillation) -- (108,-62) -| (distillation label.south);
  \node[event,anchor=north] (modular label) at (118,-68) {模块化元学习\\2018};
  \draw[leader] (modular) -- (120,-66) -| (modular label.north);
  \node[event,anchor=north] (soups label) at (140,-68) {Model Soups 2022\\KnOTS 2025};
  \draw[leader] (soups) -- (136,-66) -| (soups label.north);
\end{tikzpicture}

  \caption{三条并行研究线索中的代表工作。圆点按等距年份定位，标签通过引线对应圆点；横带区分研究问题，不规定起止时间。跨带虚线标示“任务规律”与“能力保持”两项共同问题，不表示历史继承。蒸馏取2015年arXiv版本年份；图中位置不构成“首次”判断。}
  \label{fig:kts-overview-history}
\end{figure*}

\section{本章小结}
\label{sec:kts-overview-summary}

本章为整个第三部分搭好了概念框架。跨任务利用经验有三种典型姿态：把旧任务表示拿去帮新任务的\textbf{迁移}、多个任务从头共享一套底子的\textbf{多任务学习}、以及为“以后快速学会任何新任务”做准备的\textbf{元学习}。要谈这些，先得立起来源/目标、领域/任务两组正交概念，并用风险 $R_D$ 与经验风险 $\hat R_D$ 把“在目标上好不好”变成可比较、可估计的量；协变量偏移是“只有输入边缘分布变、标签规律不变”的特例，而无标签适应在缺乏“两域联系/覆盖”依据时存在硬性的不可学习边界。

三条方法论纪律值得带走。\textbf{第一，收益只能由目标侧、用对照说话}：来源侧成绩、参数确实被改动、这一次答对了，单独都不是完整证据；“用了来源”甚至可能帮倒忙（负迁移）。\textbf{第二，名字不代表权限}：领域自适应、领域泛化、测试时适应的差别全在“目标样本何时、以何种方式参与”，必须把预测执行与按目标样本调参分开；两个都叫“迁移学习”的方案，能读到的样本、模型、反馈可以天差地别。\textbf{第三，“学会新的”和“别弄坏别的”要分别验证}：迁移要看新目标收益，共享要逐个任务看损益，融合要看来源保留与未见组合，且都要把标签、训练、推理、存储这几本不同的成本账算清。

历史上，从 Caruana 的多任务共享、Baxter 的偏置学习、Ben-David 的可学习性边界、Yosinski 的分层迁移，到学习优化器、MAML、原型网络、适配器、GPT-3 上下文与基础模型，再到集成、蒸馏、模块化组合、Model Soups 与 KnOTS，研究者把“已有经验”逐层拆成了表示、假设空间、网络层、更新规则、起点、比较方式、适配器、上下文、预测输出、模块、完整参数与参数增量等不同对象。它们长期并行、各有条件，共同支撑起后面三章：第~\ref{chap:ket-transfer-adaptation}章展开“来源经验怎样具体地帮到新目标”，第~\ref{chap:ket-sharing}章展开“多个任务怎样共同训练又互相协调”，第~\ref{chap:ket-evaluation}章则为这些收益建立统一、严谨的证据标准。
